
\section{Non-positive models and the heart-amplitude criterion}\label{sec:heart-amplitude}

Throughout this section, assume
\[
 \cK_\lambda=\Loc(\Kc_\lambda).
\]
In particular, this holds for a stably flat universal localisation: stable flatness gives
\[
 \cK_\lambda=\Loc\langle\Cone(\Sigma)\rangle,
 \qquad
 \Kc_\lambda=\thick\langle\Cone(\Sigma)\rangle.
\]

\begin{proposition}[Existence of a non-positive model]\label{prop:nonpositive-model}
The following are equivalent:
\begin{enumerate}[label=(\roman*)]
\item \(\cK_\lambda\simeq D(B)\) for a non-positive dg algebra \(B\);
\item \(\cK_\lambda\) has a compact silting generator;
\item \(\Kc_\lambda\) has a silting object.
\end{enumerate}
If \(T\) is such an object, one may take
\[
 B_T=\tau_{\le0}\RHom_R(T,T)^{\op}.
\]
\end{proposition}

\begin{proof}
If \(D(B)\simeq\cK_\lambda\) with \(B\) non-positive, the image of the free module \(B\) is compact, generates, and has no positive self-extensions.  Thus it is compact silting.

Under the standing hypothesis, the Thomason--Neeman theorem identifies the compact objects of \(\cK_\lambda\) with \(\Kc_\lambda\).  Hence a compact silting generator of \(\cK_\lambda\) belongs to \(\Kc_\lambda\) and thickly generates it.  Conversely, if \(T\) is a silting object of \(\Kc_\lambda\), then \(\Kc_\lambda=\thick(T)\), so the standing equality gives \(\cK_\lambda=\Loc(\Kc_\lambda)=\Loc(T)\).  Finally, for such \(T\), positive cohomology of \(\RHom_R(T,T)\) vanishes.  The good truncation is a quasi-isomorphism, and Keller--Morita identifies its derived category with \(\Loc(T)=\cK_\lambda\).
\end{proof}

Let \(T\) be a compact silting generator, put
\[
 B_T=\tau_{\le0}\RHom_R(T,T)^{\op},
 \qquad
 \Gamma_T=H^0(B_T),
\]
and let
\[
 F_T=T\Ltensor_{B_T}-:D(B_T)\xrightarrow{\sim}\cK_\lambda.
\]

\begin{definition}[Bounded compatibility]\label{def:r-bounded}
Assume \(\Gamma_T\) is left noetherian.  We call \(T\) \emph{\(R\)-bounded} if
\[
 F_T(\Dfg(B_T))=\Kb(\lambda).
\]
Over a field with \(\Gamma_T\) finite-dimensional, this reads
\(F_T(\Dfd(B_T))=\Kb(\lambda)\).
\end{definition}

By \Cref{cor:rbounded-generator-independent}, in finite dimension this is a property of the embedded kernel, not of the chosen silting generator.

\begin{lemma}\label{lem:d-bfg-heart}
If \(B\) is non-positive and \(H^0(B)\) is left noetherian, then
\[
 \Dfg(B)=\thick(H^0(B)\text{-}\modcat),
\]
where heart modules are viewed as dg \(B\)-modules via
\(B\to H^0(B)\).
\end{lemma}

\begin{proof}
The standard truncation triangle
\[
 \tau_{\le n-1}X\longrightarrow X\longrightarrow H^n(X)[-n]\longrightarrow
\]
reduces the cohomological amplitude.  Induction expresses every object of
\(\Dfg(B)\) as a finite extension of shifts of finitely generated heart modules.  The reverse inclusion holds because \(\Dfg(B)\) is thick.
\end{proof}

Define the \emph{heart bimodule}
\[
 M_T=T\Ltensor_{B_T}\Gamma_T.
\]
It is an \(R\)-\(\Gamma_T\) bimodule object.  For a finitely generated left \(\Gamma_T\)-module \(X\), derived associativity gives
\begin{equation}\label{eq:heart-realization}
 F_T(X)\simeq M_T\Ltensor_{\Gamma_T}X.
\end{equation}
For the general realization-functor perspective associated with hearts and tilting structures, see Psaroudakis--Vit\'{o}ria~\cite{PsaroudakisVitoria18}.

\begin{definition}[Fixed finite Tor-amplitude]\label{def:fixed-tor-amplitude}
The right \(\Gamma_T\)-complex \(M_T\) has \emph{finite Tor-amplitude} if there exist fixed integers \(a\le b\) such that, for every finitely generated left \(\Gamma_T\)-module \(X\),
\[
 M_T\Ltensor_{\Gamma_T}X\in\Db(R\text{-}\modcat)
 \quad\text{and}\quad
 H^n(M_T\Ltensor_{\Gamma_T}X)=0
 \quad(n\notin[a,b]).
\]
The endpoints are independent of \(X\).  Merely bounding the length of an interval whose position may vary with \(X\) is not sufficient and is not used here.
\end{definition}

\begin{proposition}[Noetherian direction]\label{prop:noetherian-heart}
Assume \(R\) and \(\Gamma_T\) are left noetherian.
\begin{enumerate}[label=(\alph*)]
\item If \(M_T\) has finite Tor-amplitude, then
\[
 F_T(\Dfg(B_T))\subseteq\Kb(\lambda).
\]
\item If, in addition,
\begin{equation}\label{eq:dagger}
 (\dagger_T)\qquad
 \RHom_R(T,X)\in\Dfg(B_T)
 \quad\text{for every }X\in\Kb(\lambda),
\end{equation}
then \(T\) is \(R\)-bounded.
\item If \(\Gamma_T\) is artinian and \(T\) is \(R\)-bounded, then \(M_T\) has finite Tor-amplitude.
\end{enumerate}
\end{proposition}

\begin{proof}
By \Cref{lem:d-bfg-heart}, \(\Dfg(B_T)\) is thickly generated by finitely generated \(\Gamma_T\)-modules.  Finite Tor-amplitude and \eqref{eq:heart-realization} place the image of every such heart module in \(\Kb(\lambda)\); exactness gives (a).

For (b), let \(X\in\Kb(\lambda)\).  The Morita counit is an isomorphism
\[
 F_T\RHom_R(T,X)\xrightarrow{\sim}X.
\]
Condition \((\dagger_T)\) places the inverse image in \(\Dfg(B_T)\), giving the reverse inclusion.

For (c), let \(L_1,\ldots,L_r\) be the finitely many simple \(\Gamma_T\)-modules.  Since \(T\) is \(R\)-bounded, each
\[
 M_T\Ltensor_{\Gamma_T}L_i\simeq F_T(L_i)
\]
lies in \(\Kb(\lambda)\).  Choose common lower and upper cohomological bounds \(a,b\) for this finite family.  Every finitely generated \(\Gamma_T\)-module has finite length.  Induction on length, using triangles induced by short exact sequences, gives the same fixed bounds \([a,b]\) for all such modules.
\end{proof}

\begin{theorem}[Finite-dimensional heart-amplitude criterion]\label{thm:heart-amplitude}
Let \(R\) be a finite-dimensional \(k\)-algebra, assume \(\cK_\lambda=\Loc(\Kc_\lambda)\), and let \(T\) be a silting object of \(\Kc_\lambda\), hence a compact silting generator of \(\cK_\lambda\).  Then the following are equivalent:
\begin{enumerate}[label=(\roman*)]
\item \(T\) is \(R\)-bounded;
\item \(M_T=T\Ltensor_{B_T}\Gamma_T\) has fixed finite Tor-amplitude over \(\Gamma_T\).
\end{enumerate}
When these conditions hold,
\[
 \DeltaB(\lambda)
 \simeq
 \Dfd(B_T)/\per(B_T).
\]
Moreover, by \Cref{cor:rbounded-generator-independent}, if the condition holds for one compact silting generator, it holds for every compact silting generator of the kernel.
\end{theorem}

\begin{proof}
The implication (i)\(\Rightarrow\)(ii) is
\Cref{prop:noetherian-heart}(c), since \(\Gamma_T\) is finite-dimensional.

Assume (ii).  Part (a) of \Cref{prop:noetherian-heart} gives
\[
 F_T(\Dfd(B_T))\subseteq\Kb(\lambda).
\]
Conversely, if \(X\in\Kb(\lambda)\), then \(X\) is represented by a bounded complex of finite-dimensional \(R\)-modules, while \(T\) is a bounded complex of finitely generated projectives.  Hence
\(\RHom_R(T,X)\) has finite-dimensional total cohomology, so it belongs to \(\Dfd(B_T)\).  The Morita counit identifies its image with \(X\).  Thus equality holds and \(T\) is \(R\)-bounded.

Finally, \(F_T\) sends \(\per(B_T)\) to
\(\thick(T)=\Kc_\lambda\).  Quotienting gives the formula.  Generator independence is \Cref{cor:rbounded-generator-independent}.
\end{proof}

